\chapter{Urethrotomy}

\section*{Overview}
Urethrotomy is a procedure that cuts a urethral stricture to widen the urethra and restore the flow of urine.

\section*{Alternative Names}
Direct vision internal urethrotomy; optical urethrotomy; urethral stricture surgery

\section*{Principle}
Under direct vision, an instrument with a blade cuts through the stricture, allowing the urethra to heal at a wider caliber.

\section*{History}
Optical urethrotomy was introduced in the 1970s and became the first-line endoscopic treatment for short urethral strictures.

\section*{Why This Test or Procedure Is Ordered}
Performed for urethral strictures causing obstructed voiding, urinary tract infection, or acute urinary retention.

\section*{Is This a Primary or Secondary Test or Procedure}
Primary endoscopic treatment for short strictures.

\section*{How This Test or Procedure Is Performed}
Performed under general or regional anesthesia. A urethrotome is passed to the stricture, and the blade cuts it, after which a catheter is placed.

\section*{What Preparation Is Needed}
Imaging (retrograde urethrogram) defines the stricture.

\section*{Contraindications and Precautions}
Long or dense strictures may require open urethroplasty instead.

\section*{Risks}
Recurrence, bleeding, infection, and injury to the urethra.

\section*{What Happens After the Test or Procedure}
A urinary catheter remains for days; voiding is reassessed.

\section*{Understanding Your Results}
Successful voiding and imaging confirm a widened urethra.

\section*{Normal Results}
A patent urethra with normal urinary flow.

\section*{Follow-up Tests and Procedures}
Uroflowmetry and repeat procedures if the stricture recurs.

\section*{References}
\begin{enumerate}
  \item MedlinePlus. Urethrotomy. U.S. National Library of Medicine; 2025.
  \item Current Medical Diagnosis and Treatment. McGraw-Hill; 2025.
\end{enumerate}

\clearpage
